% 【基础正文】所属：第2章（由 chapters/revised/02-skepticism-luck.tex \input 引入）
\minitopic{怀疑论不是一般的多疑}哲学怀疑论不是“什么都不信”的性格，而是针对某类知识可能性的论证。笛卡尔通过梦和恶魔设想，把可怀疑的信念逐层撤回，以寻找不可动摇的基础\parencite{descartes1996}。当代版本常以“缸中之脑”或模拟世界表达：主体的全部经验可能与正常世界相同，却没有通常以为的外部对象。

怀疑论的力量来自三个看似可信的原则。设 $O$ 为“我有两只手”，$S$ 为“我不是无手的缸中之脑”。
\begin{argumentbox}
\begin{enumerate}
  \item 如果我知道 $O$，并且知道 $O$ 蕴涵 $S$，那么我知道 $S$。（知识的闭合）
  \item 我不知道 $S$。（无法排除怀疑情景）
  \item 因此，我不知道 $O$。
\end{enumerate}
\end{argumentbox}
论证是有效的；回应者必须拒绝或限制至少一个前提。仅仅说怀疑情景“很荒唐”不够，因为问题正是：我们凭什么把它排除？

\minitopic{不可错是否是知识的标准}怀疑论往往借用了\term{不可错论}{infallibilism}：若主体可能在相同根据下出错，就没有知识。但日常和科学知识似乎都是可错的。医院根据高准确率检测作出诊断，历史学家根据残缺档案重建事件；承认未来可能修正，不等于当下没有知识。\term{可错论}{fallibilism}主张：知识要求足够好的认识位置，却不要求逻辑上排除一切错误可能。 困难在于“足够好”不能只解释为高概率。彩票中每张票几乎必输，我们通常仍不说自己知道某一张会输。反过来，普通知觉并无可计算概率，却常构成知识。相关可能性、证据种类和信念形成方式都会影响认识地位。

\minitopic{三条回应路线}\minitopic{摩尔式反转}摩尔举起双手，主张自己知道“这里有一只手，那里有另一只手”，从而知道外部世界存在\parencite[144--145]{moore1959}。其逻辑可把怀疑论倒转：我比“我不知道自己不是缸中脑”更确信自己有手，因此应拒绝后者。摩尔没有提供怀疑者愿意接受的中立证明；他的贡献是指出，怀疑论前提并不天然比日常知识更可信。

\minitopic{语义与外在主义回应}普特南论证，若一个主体一直生活在缸中，其“缸”和“脑”等词未必指称真实缸脑；思想的内容部分依赖与环境的因果联系\parencite{putnam1981}。这可能削弱某些自我指涉的怀疑假设，却不能一举保证全部日常知识：近期被放入缸中的主体仍可形成有意义的怀疑。 外在主义还可说，知识不要求主体能够从内部证明自身可靠；只要知觉事实上由合适的可靠关系产生，主体就可能知道。怀疑者会追问，这是否只是从外部替主体宣布胜利，而没有满足其反思要求。

\minitopic{相关可能性与语境}另一回应认为，知道 $p$ 只要求排除语境中相关的反例。正常情境下，模拟假设没有实际根据，因而不妨碍“我有手”；哲学课堂把它显著化后，知识声称的标准升高。此路线保留了日常断言，却必须解释语境为何改变“知道”的真值，而不仅改变说话是否谨慎。

\minitopic{标准问题与认识循环}怀疑论还有更一般的形式：判断哪些信念可靠，需要一个可靠标准；判断标准可靠，又似乎需要先知道一些信念。奇泽姆把这称为\term{标准问题}{problem of the criterion}。方法主义先提出标准，再筛选知识；特殊主义先承认若干明显知识，再提炼标准。两者都可能被指责为武断或循环。 并非所有循环都同样致命。任何总体认识方法的辩护都很难完全站在认识活动之外。关键是区分“用结论直接重述前提”的恶性循环，与多个来源相互校正、能够发现错误的开放式反思。

\begin{comparison}
《庄子》通过梦、视角转换和“彼亦一是非，此亦一是非”动摇自以为固定的区分\parencite[chap.~2]{zhuangzi2009}。它与笛卡尔怀疑都使熟悉经验陌生化，但目标并不相同：笛卡尔寻求确定基础，庄子常质疑执着于单一尺度。把两者都叫“怀疑论”只能标记问题家族，不能证明它们共享结论。
\end{comparison}

\minitopic{原典细读：怀疑的层级}笛卡尔的怀疑并非一次把所有信念清空。感官偶尔欺骗首先挑战远处或不利条件下的知觉，却不足以推翻近处清楚知觉；梦论证进一步挑战整个当前经验场景，但似乎仍保留数学和形式结构；恶魔设想才把错误可能扩展到最简单推理\parencite{descartes1996}。逐层升级很重要，因为每一步需要不同回应。用“感官会错”直接推出“没有知觉知识”，中间遗漏了从偶尔错误到从不构成知识的强前提。 休谟式怀疑更关注推理资格。我们事实上会根据经验期待未来，但理性似乎既不能演绎证明自然保持一致，也不能不循环地用过去成功证明归纳\parencite[sec.~4]{hume2000}。这里挑战的不是外部世界是否存在，而是某种推理规则如何得到规范辩护。把笛卡尔与休谟都称为怀疑论没有错，却必须说明对象分别是经验真实性还是归纳正当性。

\minitopic{怀疑论回应的代价表}任何回应都保存一些直觉并牺牲另一些。不可错论保存知识的最高标准，却使日常和科学知识大幅缩减；可错论保存这些知识，却须解释允许多大风险。拒绝闭合保存“我有手”，却使演绎扩展变得不稳定；语境主义保存日常话语与哲学怀疑的双重直觉，却承担语义证据负担；外在主义保存儿童和动物知识，却可能不能满足第一人称反思要求；摩尔式回应维护常识，却无法给怀疑者一套共同起点。 制作代价表是一种普遍方法。哲学理论很少在所有维度上占优。比较时应固定评价标准：解释直觉的范围、理论简洁性、与其他原则的兼容性，以及能否指导实际判断。仅列出一个反常结论并不足以击败理论，因为竞争理论可能付出更大代价。

\minitopic{案例工作坊：深度伪造}一段视频在视觉和声音上无法与真实录像区分。主体仅看视频便相信某人发表了言论。与缸中脑不同，深度伪造不是全局怀疑情景，而是现实中具有基础率、来源元数据和传播动机的局部风险。因此合理回应不是停止相信所有视频，而是调整相关可能性：来源不明、政治利害高、缺乏独立报道时，伪造可能性变得相关；私人实时通话且有交叉验证时，标准不同。 这个例子说明，怀疑情景一旦进入现实，就从“逻辑上不能排除”转化为证据管理问题。主体可检查原始文件、时间线、其他视角和发布者记录。可错知识不要求消灭一切伪造可能，但要求风险与核验强度相称。怀疑论训练的积极意义，是帮助我们识别何时应提高标准，而不是培养无差别不信任。
